\subsubsection{Auction and Dealer Market Structures}

\begin{taxonomy} \hlt{Execution-System Segmentation}
\begin{enumerate}[label=\roman*.]
\setlength{\itemsep}{0pt}
\item \hlt{quote-driven}: dealers supply liquidity through tradable quotes and balance sheets;
\item \hlt{order-driven}: rules match public orders under priority and pricing rules;
\item \hlt{brokered}: brokers search for counterparties and reveal interest selectively;
\item \hlt{hybrid}: venue rules combine books, dealers, auctions, and negotiation.
\end{enumerate}
\end{taxonomy}

\begin{taxonomy} \hlt{Trading Session Segmentation}
\begin{enumerate}[label=\roman*.]
\setlength{\itemsep}{0pt}
\item \hlt{continuous}: trades can occur whenever compatible orders arrive;
\item \hlt{call}: orders batch until a scheduled or triggered uncrossing;
\item \hlt{rotation}: instruments are called sequentially;
\item \hlt{extended}: trading occurs outside the regular session, usually with thinner liquidity.
\end{enumerate}
\end{taxonomy}

\begin{mechanism} \hlt{Open-Outcry Oral Auctions}\\
Open-outcry markets arrange trades through public bids, offers, and acceptances on a floor or pit. The economic idea is still useful electronically: trading interest should be exposed enough for competing traders to react, while the rule book decides whose interest has precedence.
\end{mechanism}

\begin{marketrule} \hlt{Public Exposure Rule}\\
An auction is open when bids, offers, and acceptances are announced to the crowd before a trade is final. Public exposure supports fairness and price competition, but it also reveals trading interest.
\end{marketrule}

\begin{taxonomy} \hlt{Order-Driven Rule Stack}
\begin{enumerate}[label=\roman*.]
\setlength{\itemsep}{0pt}
\item \hlt{precedence rules}: rank orders when more than one trader wants the same trade;
\item \hlt{matching rules}: decide which buy and sell orders can trade;
\item \hlt{pricing rules}: determine the transaction price once a match is feasible;
\item \hlt{access rules}: decide who may submit, observe, cancel, or route orders.
\end{enumerate}
\end{taxonomy}

\begin{mechanism} \hlt{Continuous Auctions}\\
Orders execute whenever they arrive and match resting interest. Continuous trading provides immediacy but exposes prices to asynchronous order-flow imbalances.
\end{mechanism}

\begin{mechanism} \hlt{Call Auctions}\\
Orders accumulate, then execute at a single clearing price. Opening and closing calls concentrate liquidity when uncertainty and imbalance are high.
\end{mechanism}

\begin{mechanism} \hlt{Fixing Auctions}\\
Fixing auctions are call auctions used to create reference prices.\\
They are valuable when investors need one public price for index valuation, NAV computation, derivative settlement, opening risk transfer, or end-of-day inventory adjustment.
\end{mechanism}

\begin{principle} \hlt{Reference Price Monopoly}\\
Even when continuous trading fragments, primary exchanges often retain the auction flow that defines official reference prices. This makes auction volume structurally different from continuous-session volume and gives the primary market a separate source of liquidity concentration.
\end{principle}

\begin{taxonomy} \hlt{Auction Session Segmentation}
\begin{enumerate}[label=\roman*.]
\setlength{\itemsep}{0pt}
\item \hlt{opening call}: rebuilds the book after overnight information and sets the first tradable price;
\item \hlt{intraday call}: concentrates liquidity around scheduled events or low-liquidity periods;
\item \hlt{closing call}: sets the end-of-day reference price and absorbs benchmark-driven flow;
\item \hlt{volatility call}: pauses continuous trading when price moves or imbalances breach venue rules.
\end{enumerate}
\end{taxonomy}

\begin{mechanism} \hlt{Dealer Markets}\\
Prices form through dealer quotes rather than a single public book. Dealers absorb timing gaps, manage inventory, and revise quotes as information and inventory change.
\end{mechanism}

\begin{mechanism} \hlt{Brokered Markets}\\
Brokers create liquidity by search rather than inventory commitment. This fits blocks, unique assets, and illiquid instruments where public standing orders would reveal too much or attract too little interest.
\end{mechanism}

\begin{taxonomy} \hlt{Hybrid Markets}\\
Hybrid markets combine order books, dealers, auctions, hidden orders, and off-book negotiation. They mix public order interaction with intermediary support.
\end{taxonomy}

\begin{marketrule} \hlt{Priority and Access Rules}\\
Priority, tick size, order type, display, halt, and access rules determine how trades are resolved. The same latent supply and demand can produce different prices, quantities, and costs under different rules.
\end{marketrule}

\begin{casestudy} \hlt{Gold Fixing as Sequential Auction}\\
The London gold fixing resembles a call process: participants submit net interest, revise after tentative prices, and stop when net demand and supply balance. It produces a clearing price, not continuous immediacy.
\end{casestudy}

\begin{practitionernote} \hlt{Why Large Orders Like Calls}\\
Calls reduce the need to reveal urgency sequentially. A large order can interact with accumulated opposite interest at one price, often with less signalling than chasing liquidity through the continuous book.
\end{practitionernote}

\begin{caveat} \hlt{Auction Transparency Trade-Off}\\
Pre-auction imbalance messages help traders converge to a clearing price, but they can also invite probing, cancellation, and last-second repricing. More transparency improves price discovery only when participants cannot cheaply game the signal.
\end{caveat}
